\documentclass[conference]{IEEEtran}
\usepackage[utf8]{inputenc}
\usepackage[T1]{fontenc}
\usepackage{cite}
\usepackage{amsmath,amssymb}
\usepackage{graphicx}
\usepackage{booktabs}
\usepackage{url}

\title{Measuring Validator and Sandboxed-Execution Coverage for LLM‑Generated NetOps Artifacts: A Paired Confirmatory Study}

\author{
\IEEEauthorblockN{Autonomous AI Researcher}
\IEEEauthorblockA{CNSM 2026 Agentic AI Researcher Track}
}

\begin{document}

\maketitle

\begin{abstract}
We preregistered and executed a paired confirmatory experiment (study\_id f2c3d8b1-7a6e-4a9a-b2c3-5f8e1d2b6c7e) that measured the marginal detection benefit of adding sandboxed emulation to a deterministic validator for LLM-generated intent\ensuremath{\rightarrow}configuration artifacts. Under shared\_initial\_candidate semantics using the declared model "gpt-5-mini" (provider: openai\_responses) we evaluated Npairs = 720 confirmatory tasks. The deterministic validator (deterministic\_netops\_validator\_v5) flagged 719/720 = 0.9986111111111111; the guarded pipeline (validator + sandboxed emulation) flagged 720/720 = 1.0. Paired difference (guarded - baseline) = 0.001388888888888884; exact McNemar two-sided p = 1.0; 95\% bootstrap percentile CI (10000 resamples, seed = 1729) = [0.0, 0.004166666666666667]. Run accounting: model\_calls\_used = 721 (720 shared initial + 1 repair-stage call). These artifact-grounded results lie far below our preregistered minimum effect-of-interest (0.20); we report the preregistered design, exact quantitative outcomes, run-accounting diagnostics, artifact-grounded interpretation for the negligible guarded benefit in this regime, and reproducibility pointers into the archived execution bundle.
\end{abstract}

\section{Introduction}
Generative large language models (LLMs) are increasingly proposed to translate operator intent into device configuration and NetOps workflows. Operational deployment requires generated artifacts to satisfy syntactic correctness, workflow ordering, and higher-order safety/policy constraints. A common mitigation architecture is generate\ensuremath{\rightarrow}validate\ensuremath{\rightarrow}(repair)\ensuremath{\rightarrow}execute, sometimes augmented by sandboxed emulation to reveal runtime-only failures that static validators miss. Prior surveys and system reports advocate such workflow-centred mitigations [1]--[3], but provide limited large-scale paired measures of the marginal value added by sandboxed execution. Motivated by this gap, we preregistered and executed a controlled paired study to measure the aggregate detection-rate difference contributed uniquely by an automated sandboxed-emulation stage when the same initial generated candidate is analyzed by both conditions.

Research question and contributions. Our preregistered research question was: for synthetic intent\ensuremath{\rightarrow}configuration tasks under shared\_initial\_candidate semantics, what is the paired detection-rate difference (guarded - baseline) between (A) a deterministic validator-only pipeline and (B) the same deterministic validator followed by automated sandboxed emulation? Contributions: (1) a machine-readable preregistration and faithful autonomous execution; (2) an artifact-grounded paired analysis on Npairs = 720 with complete accounting; (3) an artifact-grounded interpretation explaining a near-zero marginal effect; (4) concise reproducibility pointers into the archived bundle and prescriptive boundary conditions where guarded stages are likely to matter.

\section{Related Work}
Workflow-centred mitigations and verification tools. Surveys and recent system reports document architectures that pair LLM generation with validators, verifiers, and emulators to improve correctness of generated NetOps artifacts and adjacent program-generation tasks [1], [2]. Tool-focused work demonstrates generate\ensuremath{\rightarrow}validate\ensuremath{\rightarrow}repair patterns (e.g., Batfish/NetKAT style verification, emulator-assisted validation) but reports heterogeneous gains depending on task, validator rule coverage, and emulator fidelity [2], [3].

Empirical findings motivating a paired design. Empirical studies in program/specification generation show external verifiers often catch model mistakes not found by the model alone, and some NetOps prototypes use emulation to expose runtime issues [3]. These prior findings motivated our preregistered paired design that isolates guarded-stage contributions: because we enforce shared\_initial\_candidate semantics, the only permissible sources of between-condition discordance are downstream guarded-stage emulation observations and permitted repair-stage invocations recorded in run accounting.

\section{Methodology}
Preregistration and experimental design. We preregistered a confirmatory paired-binary design (study\_id f2c3d8b1-7a6e-4a9a-b2c3-5f8e1d2b6c7e). Primary estimand: paired\_success\_rate\_difference\_guarded\_minus\_baseline aggregated across confirmatory samples. Planned confirmatory sample: Npairs = 720 (planned episodes = 1440). Minimum effect-of-interest for power planning = 0.20 (20 percentage points). The preregistration specified inclusion/exclusion rules, missingness handling, multiplicity handling for exploratory secondary tests, and deterministic contamination controls.

Task bank, vendors, and inclusion. Confirmatory items were synthetic intent\ensuremath{\rightarrow}configuration tasks from intent\_configuration\_repair\_v1 covering three abstracted vendor-CLI families and three topology templates (leaf-spine, 3-node service chain, triangle). Inclusion required vendor-specific parser acceptance; the preregistered missingness policy permitted up to two deterministic reseeds for parser failures; irrecoverable parser failures would be deterministically excluded and replaced. In execution there were zero irrecoverable parser exclusions: complete\_pair\_count = 720 (see analysis/missingness\_summary.csv).

Model configuration and sampling semantics. Declared/requested model: "gpt-5-mini" (provider field: openai\_responses). The archived model configuration records max\_output\_tokens = 2000 and temperature = null (execution/model\_configuration.json, sha256 a40e96e212ab87da...). Important clarification (preregistered and enforced): temperature = null indicates the client did not set an explicit temperature parameter; null does not imply temperature=0 or provider-side determinism. We enforced shared\_initial\_candidate semantics: for each pair we obtained exactly one sampled initial candidate (shared) which both baseline and guarded pipelines analyzed. Therefore any between-condition difference can only arise from (i) guarded-stage emulation observations not encoded as a validator rule, or (ii) the allowed downstream repair-stage invocation(s) recorded in run accounting.

Validator and guarded pipeline implementation. Baseline: deterministic\_netops\_validator\_v5 performing full intent-constraint validation (syntax, command ordering/workflow, required-field presence, and policy-level constraints) and returning binary detected/not-detected per the preregistered scoring rule full\_intent\_constraint\_validation. Guarded: baseline validator followed by an automated sandboxed-emulation harness executing the candidate against an emulated instance of the declared topology and producing runtime observations. The repair policy permitted at most one repair-stage model call per task; run accounting records exactly one repair-stage call in the confirmatory execution.

Execution controls, provenance, and deterministic reconciliation. The adapter family used: hosted\_netops\_gvr\_v1. Execution manifest and deterministic reconciliation artifacts record completed\_episode\_count = 1440, complete\_pair\_count = 720, failed\_episode\_count = 0, and provider-call audit information including model\_calls\_used = 721 and stage\_counts: shared\_initial\_generation = 720; repair = 1. Provenance and contamination controls (deterministic hashing, artifact-version checks) were applied per the preregistration; no contamination flags occurred for the confirmatory set.

\section{Results}
Primary confirmatory counts and test statistics (artifact-grounded). Analysis artifacts (analysis/paired\_contingency\_table.csv and analysis/condition\_summary.csv) report the contingency counts: n11 = 719 (detected by both), n10 = 1 (guarded-only), n01 = 0 (baseline-only), n00 = 0 (neither). Marginal detection rates: baseline = 719/720 = 0.9986111111111111; guarded = 720/720 = 1.0. Paired estimate (guarded - baseline) = 0.001388888888888884. Exact McNemar two-sided test p = 1.0. Paired-bootstrap percentile 95\% CI (10000 resamples, bootstrap\_seed = 1729) = [0.0, 0.004166666666666667]. All numeric values below are reproduced verbatim from analysis/results.json and analysis artifacts.

Table 1 --- Primary confirmatory summary (artifact-grounded).
- Npairs = 720 (complete paired episodes)
- Baseline successes = 719; Baseline success rate = 0.9986111111111111
- Guarded successes = 720; Guarded success rate = 1.0
- Paired difference (guarded - baseline) = 0.001388888888888884
- Exact McNemar two-sided p = 1.0
- 95\% bootstrap percentile CI = [0.0, 0.004166666666666667] (10000 resamples, seed = 1729)

Discordant-pair attribution and repair accounting. Provider-call accounting and deterministic reconciliation identify a single repair-stage invocation (stage\_counts: repair = 1) and a unique n10 discordant count. Because shared\_initial\_candidate semantics were enforced, the permissible causes for that single guarded-only detection are: (i) a guarded-stage emulation observation not encoded as a validator violation, or (ii) the single recorded repair-stage invocation. The confirmatory execution bundle contains the raw per-episode records; the unique discordant record can be located by searching execution/raw\_results.jsonl (input-results SHA-256 = df367ecc\allowbreak{}40a6f0ea\allowbreak{}f4021292\allowbreak{}60e11d84\allowbreak{}50884a4e\allowbreak{}a5fdb94a\allowbreak{}602e0d25\allowbreak{}13aef0e8) for the single entry where baseline\_score = 0 and guarded\_score = 1. Deterministic reconciliation is archived as analysis/deterministic\_reconciliation.json (sha256 49731a2e\allowbreak{}52e049cc\allowbreak{}f5b53c35\allowbreak{}ddac2351\allowbreak{}ad75e7e9\allowbreak{}5504dc4a\allowbreak{}42bcb096\allowbreak{}d0dec3aa) and its provider\_call\_audit confirms model\_calls\_used = 721 and repair = 1. The raw\_results record for the discordant entry contains the pair\_id and the validator\_trace and repair\_provider\_trace fields that link the guarded detection to the documented repair-stage call; these artifacts are part of the canonical execution bundle referenced by the execution manifest.

Representative per-episode diagnostics. Representative scoring traces archived under execution/scoring/task-000001-*.json ... task-000006-*.json include validator\_trace objects with normalized\_configuration, validation\_before/after, parsed\_command\_count, required\_command\_count, violation\_count, and validator\_version = 5.0. Those representative artifacts illustrate that the vast majority of shared\_initial candidates complied with validator rules (violation\_count = 0) and that sandboxed emulation largely corroborated validator findings in this curated task bank.

\section{Discussion}
Why was the guarded-stage aggregate benefit negligible? The archived evidence supports three artifact-grounded factors and suggests concrete boundary conditions where guarded stages remain important.

1) Validator ceiling effect (artifact-grounded). The deterministic validator covered the explicit syntactic and semantic constraints exercised by the synthetic tasks: most shared\_initial candidates parsed successfully and produced validator traces with violation\_count = 0. Analysis artifacts record n11 = 719 and n10 = 1, so the baseline validator already detected or accepted nearly all conformant configurations (baseline\_success\_rate \ensuremath{\approx} 0.9986). With validator coverage at this level there is little statistical headroom for the guarded stage to contribute aggregate detections at scale. Representative validator\_trace fields (parsed\_command\_count and required\_command\_count in execution/scoring/*.json) confirm that tasks were largely within the validator's rule set.

2) Task curation and inclusion policy concentrate cases inside validator coverage. The confirmatory bank used parser-validated synthetic items that explicitly encoded workflow patterns (shutdown\ensuremath{\rightarrow}configure\ensuremath{\rightarrow}restore, make-before-break uplink switchover, paired MTU changes). The inclusion rule required vendor-parser acceptance before pairing, and the missingness artifact shows zero irrecoverable parser failures for the confirmatory set (analysis/missingness\_summary.csv). This curation intentionally constrained the evaluation to cases where the validator's deterministic rules apply; it therefore reduces the prevalence of subtle runtime-only failure modes (race conditions, unexpected protocol interactions, or stateful side-effects) that emulation would be uniquely positioned to reveal.

3) Shared\_initial\_candidate semantics and observed repair activity. Per the preregistration we enforced shared\_initial\_candidate: the same single sampled candidate was analyzed by both conditions for each pair. Provider-call accounting and deterministic\_reconciliation.json show model\_calls\_used = 721 (720 shared initial + 1 repair-stage call) and stage\_counts repair = 1. Under these semantics, between-condition discordance cannot arise from different model outputs; it can only arise from (a) additional runtime observations surfaced by emulation that are not encoded as validator rules, or (b) permitted downstream repair actions. The single n10 discordant count is explicitly mapped in the archived raw\_results.jsonl record to the recorded repair-stage invocation, so the observed discordance is attributable to downstream guarded-stage activity rather than generation variability.

Practical interpretation and policy implications. The artifact-grounded conclusion is conditional: when a high-coverage static validator exists and inputs are constrained and parser-validated, automated sandboxed emulation may add negligible marginal detection at scale. Operationally, this suggests a cost/benefit trade-off: expensive emulation and canary infrastructure will have the most value when (i) validator coverage is known to be incomplete for target workflows, (ii) inputs include unconstrained or noisy natural-language intents, (iii) stateful multi-step changes introduce emergent runtime behaviors, or (iv) adversarial or fuzzed inputs are expected. In those regimes emulation can expose violations of invariants not expressible in the validator's static rule set (for example, transient traffic blackholing under failure-induced routing changes or vendor-specific CLI side-effects). The archived per-episode traces (execution/scoring/*.json) provide concrete examples and can be used to build such adversarial testbeds.

Diagnostics useful to operators and evaluators. The run artifacts enable targeted diagnostics rather than blanket conclusions. For example:
- Coverage auditing: compare validator\_rule\_inventory to parsed\_command\_count statistics across tasks to locate missing rules; analysis/condition\_summary.csv and execution/scoring/* contain the necessary fields.
- Discordant attribution: locate the unique discordant raw\_results.jsonl entry (baseline\_score = 0 and guarded\_score = 1) and follow its validator\_trace and repair\_provider\_trace to see whether the guarded detection arose from an emulation observation or an applied repair; deterministic\_reconciliation.json confirms the single repair-stage call linkage.
- Emulation fidelity checks: examine the emulation observation summaries in validator\_trace.validation\_after and compare them to expected\_final\_state in task\_manifest entries to identify classes of runtime behaviors the emulator notices but the validator cannot express.
These diagnostics are artifact-grounded and reproducible from the archived bundle referenced in the execution manifest.

Threats to validity and scope (artifact-grounded). Internal: shared\_initial\_candidate reduces variation in generation and isolates guarded-stage effects but constrains external generality --- in practice operators may allow multiple candidate generations or stochastic sampling to improve correctness, which could change the relative benefit of guarded stages. Construct: our tasks are synthetic and parser-validated; they target representative NetOps patterns but do not exhaust the space of runtime faults (e.g., vendor implementation bugs, timing-sensitive state races). External: a single declared/requested model family (gpt-5-mini) and one validator implementation (deterministic\_netops\_validator\_v5) limit generalization across models, validator families, and larger production topologies. Conclusion: power planning targeted a 0.20 minimum detectable paired difference; the observed paired difference (\textasciitilde{}0.0014) lies far below that sensitivity and therefore should be interpreted as a narrowly scoped near-zero finding for this experimental regime rather than a general negative result for all NetOps workflows. All of these limitations are documented in the preregistration and in analysis artifacts (missingness and reconciliation files).

Operational recommendations and next steps grounded in the archived evidence. To meaningfully evaluate guarded-stage benefits where they matter operationally, future controlled studies (using the same artifact-oriented framework) should: (1) include adversarially generated tasks and fuzzed intents designed to evade validator rules; (2) broaden vendor-CLI families and emulate vendor-specific OS behaviors with higher-fidelity toolchains; (3) allow multiple sampled initial candidates per task to measure interplay between generation variability and guarded-stage detection; (4) instrument canary/rollback-aware scoring that penalizes dangerous but syntactically valid changes; and (5) measure cost and latency trade-offs for emulation at scale. The archived execution artifacts and per-episode traces supplied by this run provide a reproducible starting point and concrete templates for these extensions (see execution/raw\_results.jsonl and analysis/deterministic\_reconciliation.json).

Summary. The evidence archived with this run shows that, in a curated parser-validated synthetic task bank and under shared\_initial\_candidate semantics, a very strong deterministic validator already achieves near-complete coverage and the guarded sandboxed-emulation stage produced only a single additional detection tied to a documented repair-stage call. This narrow, artifact-grounded finding highlights circumstances in which validator-centric approaches may suffice, and it precisely identifies the empirical and methodological axes (task diversity, adversarial inputs, emulator fidelity, generation variability) where guarded stages are likely to provide substantive operational value.

\section{Limitations}
\begin{itemize}
\item Synthetic task bank and deterministic task generator produce limited ecological diversity and create a validator-ceiling effect; results may not generalize to noisier, adversarial, or production distributions.
\item Single declared/requested model family (gpt-5-mini) and single validator implementation (deterministic\_netops\_validator\_v5) limit generalizability across models and validator families.
\item Preregistered shared\_initial\_candidate semantics (one initial generation reused across paired conditions) isolate guarded-stage contribution but remove between-condition generation variability; this design choice constrains discordance sources to downstream guarded-stage observations or repair calls.
\item Sandboxed-emulation fidelity relative to vendor/OS production behaviors and large-scale stateful interactions is limited by the emulation harness used; higher-fidelity emulators could change outcomes in other regimes.
\item Study power planning targeted a minimum detectable paired difference of 0.20; the observed aggregate effect (\textasciitilde{}0.0014) lies far below that design sensitivity and should be interpreted as a narrowly scoped near-zero finding for this experimental regime rather than a general negative result for all NetOps workflows.
\item Provider-side nondeterminism and hidden caching remain residual uncertainties despite deterministic reconciliation procedures; these are documented in analysis/deterministic\_reconciliation.json and provider\_call\_audit artifacts.
\end{itemize}

\section{Conclusion}
We preregistered and executed a paired confirmatory study comparing a strong deterministic validator with the same validator augmented by sandboxed emulation on a curated synthetic intent\ensuremath{\rightarrow}configuration task bank. Over 720 paired tasks the guarded pipeline produced a negligible aggregate detection gain (0.001388888888888884) with exact McNemar p = 1.0 and 95\% bootstrap CI [0.0, 0.004166666666666667]. Run accounting shows one repair-stage invocation corresponding to the single guarded-only detection; because shared\_initial\_candidate semantics were enforced, archived per-episode traces permit independent verification of that mapping via execution/raw\_results.jsonl and analysis/deterministic\_reconciliation.json. We interpret the result as arising from validator ceiling effects and curated task design; follow-up work should focus on adversarial and production-like task banks, multi-validator ensembles, and higher-fidelity emulation to identify regimes where guarded stages add substantive operational value.

Reproducibility pointers (compact). The canonical execution bundle archived by the run contains: execution/execution\_manifest.json (execution summary and preregistration SHA-256), execution/raw\_results.jsonl (input-results SHA-256 df367ecc\allowbreak{}40a6f0ea\allowbreak{}f4021292\allowbreak{}60e11d84\allowbreak{}50884a4e\allowbreak{}a5fdb94a\allowbreak{}602e0d25\allowbreak{}13aef0e8), analysis/deterministic\_reconciliation.json (sha256 49731a2e\allowbreak{}52e049cc\allowbreak{}f5b53c35\allowbreak{}ddac2351\allowbreak{}ad75e7e9\allowbreak{}5504dc4a\allowbreak{}42bcb096\allowbreak{}d0dec3aa), execution/model\_configuration.json (sha256 a40e96e2\allowbreak{}12ab87da\allowbreak{}251ac0d6\allowbreak{}dbb75bc5\allowbreak{}43afa134\allowbreak{}38b5a6b9\allowbreak{}ebd07359\allowbreak{}7ffdd820), analysis/paired\_contingency\_table.csv (sha256 efe6e8e7\allowbreak{}016fa843\allowbreak{}a8226e91\allowbreak{}1dbde829\allowbreak{}e943903d\allowbreak{}719101da\allowbreak{}8d956b3f\allowbreak{}f899133c), and analysis/condition\_summary.csv (sha256 394bc690\allowbreak{}671b7311\allowbreak{}94f9b42b\allowbreak{}4dcbe28f\allowbreak{}ded6e2ec\allowbreak{}8cf898b7\allowbreak{}98312253\allowbreak{}43c22c28). The discordant episode is the unique raw\_results.jsonl entry with baseline\_score = 0 and guarded\_score = 1; that record contains the pair\_id and validator/emulation traces that map it to the single repair-stage invocation recorded in deterministic\_reconciliation.json. These artifacts are archived as the canonical reproducibility bundle referenced by the execution manifest and the preregistration record. (See Disclosure for exact manifest references.)

\section*{Disclosure Statement}
Disclosure: Autonomous post-lock research run executed under the frozen master prompt supplied to the system. Declared/requested model: "gpt-5-mini" (provider recorded as openai\_responses). Deterministic validator: deterministic\_netops\_validator\_v5. Orchestration adapter family: hosted\_netops\_gvr\_v1. Preregistration id: f2c3d8b1-7a6e-4a9a-b2c3-5f8e1d2b6c7e (preregistration SHA-256: 0169919f\allowbreak{}b8c5aedb\allowbreak{}2c5a78e0\allowbreak{}31f1ab5c\allowbreak{}3aeda405\allowbreak{}cde1ae80\allowbreak{}fb2f8647\allowbreak{}68595c1b; recorded in the execution manifest). Initial master-prompt immutable reference: provenance/master\_prompt.txt, SHA-256 = 1872df1e\allowbreak{}1805d2d9\allowbreak{}6940456c\allowbreak{}a016bd66\allowbreak{}5d1d5196\allowbreak{}add77f5a\allowbreak{}cdf1582b\allowbreak{}b39b15ba. Execution summary (artifact-grounded): completed\_episode\_count = 1440, complete\_pair\_count = 720, model\_calls\_used = 721 (720 shared initial + 1 repair-stage call); stage\_counts and provider-call audit are archived in analysis/deterministic\_reconciliation.json (sha256 49731a2e52e0...). Human role and autonomy boundary: a human Research Director selected the broad NetOps topic family and constrained the pre-lock master prompt; after the autonomy lock the agentic pipeline autonomously performed literature selection, preregistration, code and prompt generation, experiment execution, analysis, peer review, and manuscript drafting. No manual human editing of technical manuscript content occurred after the autonomy lock. The execution manifest and archived artifacts named above serve as the canonical reproducibility bundle.

\begin{thebibliography}{99}
\bibitem{ref1} Muhammad Bilal, Jon Crowcroft, Ruizhi Wang, Xiaolong Xu, Schahram Dustdar, "Large Language Models for Agentic NetOps and AIOps: Architectures, Evaluation, and Safety", 2026, 10.48550/arxiv.2605.12729
\bibitem{ref2} Matt Brown, Ari Fogel, Daniel Halperin, Victor Heorhiadi, Ratul Mahajan, Todd Millstein, "Lessons from the evolution of the Batfish configuration analysis tool", 2023, 10.1145/3603269.3604866
\bibitem{ref3} http://arxiv.org/abs/2407.08249
\end{thebibliography}

\end{document}
